\chapter{Che cosa fa un ugello e che cosa significa ``ottimale''}
\label{ch:nozzle}

Questo capitolo non presuppone alcuna conoscenza pregressa. Spiega
qual è la macchina in esame, a che cosa serve, quali numeri la
descrivono, perché ci sia qualcosa da ottimizzare e che cosa renda i
motori a detonazione rotante un problema di progetto più difficile
dei razzi ordinari. Chi già progetta ugelli può passare al
Capitolo~\ref{ch:moc}; tutti gli altri dovrebbero leggerlo, perché
ogni capitolo successivo usa questo vocabolario senza rispiegarlo.

\section{La macchina}

Un razzo brucia propellente in una camera, producendo gas ad alta
pressione e alta temperatura ma senza una direzione privilegiata.
L'ugello è la parte che converte quella pressione immagazzinata in
uno scarico veloce e orientato --- e, per il terzo principio di
Newton, la reazione al lancio all'indietro di quello scarico è la
spinta.

La forma che fa questo è l'ugello \emph{convergente--divergente}
della Figura~\ref{fig:nozzleprimer}a: il condotto si restringe fino
a un'area minima, la \emph{gola}, e poi torna ad allargarsi. Il
restringimento accelera il gas esattamente fino alla velocità locale
del suono in gola; oltre la gola, in flusso supersonico, il
comportamento si inverte --- allargare il condotto accelera
ulteriormente il gas. Così la parte divergente, che è ciò che questo
documento progetta, continua a convertire pressione in velocità fino
all'uscita.

Tre conseguenze contano in tutto il seguito.

\emph{La gola fissa la portata.} Quando il flusso raggiunge la
velocità del suono in gola, il condotto è \emph{bloccato}: la massa
di gas al secondo, $\dot m$, è fissata dall'area di gola e dallo
stato di camera, e nulla a valle può cambiarla. Tenere fisso il
raggio di gola $y_t$ è dunque la stessa cosa che tenere fissa la
portata --- fatto che questo documento usa come vincolo di progetto.

\emph{Lo stato di camera si esprime come stato di ristagno.} La
pressione di camera $P_0$ e la temperatura $T_0$ (i valori di
\emph{ristagno} --- ciò che il gas avrebbe se portato a riposo) sono
i dati al contorno dell'intero problema. Tutto ciò che sta a valle è
lo stesso gas, espanso.

\emph{L'uscita è un compromesso con il mondo esterno.} Il gas esce a
una certa pressione $p_e$, in un'atmosfera a pressione $p_a$. Non è
detto che siano uguali, e proprio il fatto che lo siano o no è ciò
che rende il progetto non banale.

\begin{figure}[htbp]
\centering
\includegraphics[width=\textwidth]{fig_nozzle_primer.pdf}
\caption{La macchina, ciò che fa e perché ha un ottimo. (a)~Un
ugello convergente--divergente: il gas dalla camera (a sinistra)
accelera fino a condizioni soniche in gola e poi supersoniche nella
parte divergente; gli assi sono posizione assiale $x$ [m] e raggio
$y$ [m], con raggio di gola $y_t = 1$~m in questo caso di
riferimento. (b)~Che cosa fa davvero l'ugello al gas, misurato con
il solutore di questo documento lungo la parete: il numero di Mach
(blu, asse sinistro, adimensionale) cresce mentre la pressione
statica (verde, asse destro, bar, logaritmica) cala di due ordini di
grandezza --- si sta scambiando pressione con velocità. (c)~Perché
esista un ottimo: la pressione di uscita $p_e$ [bar] cala al
crescere del rapporto di espansione $\varepsilon$ [--]. Sopra la
linea dell'ambiente conviene ancora allungare l'ugello; sotto, la
lunghezza in più costa spinta. La stella segna il punto in cui si
incrociano --- l'ottimo, che il Capitolo~\ref{ch:firstopt} trova
\emph{alla cieca} e verifica contro questo incrocio.}
\label{fig:nozzleprimer}
\end{figure}

\section{I numeri che descrivono un progetto}

Bastano poche grandezze per enunciare un problema di progetto, e
questo documento si attiene a quelle (l'elenco completo dei simboli
è la Tabella~\ref{tab:notation}).

Il \emph{rapporto di espansione} $\varepsilon = A_e/A_t$ --- area di
uscita diviso area di gola --- è il numero singolo più importante:
dice quanto il gas viene espanso. Un $\varepsilon$ maggiore significa
un ugello più lungo e più largo, uno scarico più veloce e più
freddo, una pressione di uscita più bassa. Il \emph{contorno di
parete} è la forma stessa, la curva $y(x)$ dalla gola all'uscita. La
\emph{lunghezza} $L$ è il punto in cui quella curva si ferma. E
$p_a$, la pressione ambiente, è l'ambiente: circa 1~bar a livello
del mare, essenzialmente nulla nel vuoto.

\section{La spinta, e perché esiste un ottimo}

La spinta è un bilancio di quantità di moto. Scritta sul piano di
uscita è
\begin{equation}
J \;=\; \dot m\,u_e \;+\; (p_e - p_a)\,A_e ,
\label{eq:thrustintro}
\end{equation}
dove il primo termine è la quantità di moto portata fuori dallo
scarico (portata per velocità di uscita) e il secondo la forza netta
di pressione sull'area di uscita. Il Capitolo~\ref{ch:score} calcola
questa quantità in due modi indipendenti come autoverifica, ma la
fisica dell'ottimo è già visibile nella~\eqref{eq:thrustintro}.

Si supponga di allungare leggermente l'ugello, aggiungendo un
sottile anello di parete. Il gas interno spinge su quell'anello con
pressione $p$; l'atmosfera spinge in senso opposto con $p_a$. Il
guadagno netto è proporzionale a $p - p_a$. Dunque, finché lo
scarico è a pressione maggiore dell'ambiente --- il flusso è
\emph{sotto-espanso} --- la lunghezza in più paga. Una volta che il
flusso si è espanso oltre l'ambiente --- \emph{sovra-espanso} ---
ogni anello aggiuntivo spinge dalla parte sbagliata e costa spinta.
Il progetto migliore sta esattamente al cambio di segno:
\begin{equation}
p_e(\varepsilon^{*}) \;=\; p_a .
\label{eq:matchintro}
\end{equation}
La Figura~\ref{fig:nozzleprimer}c è quell'affermazione, misurata.
Questa singola idea --- estendere la parete finché la pressione
locale supera l'ambiente, fermarsi quando non lo fa --- ricompare in
tutto il documento in forma via via più generale: come condizione
terminale al labbro dell'ugello (Capitolo~\ref{ch:score}), come
condizione mediata su un ciclo (Capitolo~\ref{ch:cycle}) e come
regola che decide dove troncare un ugello a spina
(Capitolo~\ref{ch:plugcycle}).

Un caso particolare importante: nel vuoto $p_a = 0$, quindi
$p - p_a$ è positivo ovunque e \emph{più lungo è sempre meglio}. Non
esiste un ottimo interno, solo la lunghezza che il veicolo riesce a
portare. Qualunque ottimizzatore degno di fiducia deve riprodurre
quel comportamento, e diversi test di questo documento verificano
esattamente questo.

\section{Tre famiglie di ugello, e perché la terza è difficile}

\emph{L'ugello ideale (a lunghezza minima)} è la forma da manuale:
il suo contorno è tracciato in modo che lo scarico esca
perfettamente uniforme e parallelo. È facile da caratterizzare e qui
serve da riferimento; lo genera il Mattone~1.

\emph{Il contorno ottimizzato in spinta}, secondo la soluzione di
Rao del 1958~\cite{rao}, rinuncia all'uscita perfettamente uniforme
e pone una domanda migliore: fra tutte le pareti di una \emph{data
lunghezza}, quale produce la massima spinta? La risposta è una forma
più corta e più piena, ed è quella che i razzi reali usano. Rao la
risolse con il calcolo delle variazioni --- la matematica
dell'ottimizzazione su un'intera curva anziché su pochi numeri --- e
il Capitolo~\ref{ch:referee} confronta le sue condizioni derivate a
mano con la macchina automatica di questo programma proprio su un
contorno di quel tipo.

\emph{L'ugello a spina (o aerospike)} rovescia la geometria
(Figura~\ref{fig:bellvsplug}). Invece che il gas sia racchiuso da
una parete, la superficie solida --- la \emph{spina} --- sta sotto,
e la frontiera esterna dello scarico non è altro che l'atmosfera
circostante. Quella frontiera libera è ciò che rende interessante il
plug: non è fissata dall'hardware, quindi il pennacchio si
riconfigura quando la pressione ambiente cambia. Una campana è
ottimale a una quota; un plug è all'incirca giusto su un
intervallo, perché lo scarico finisce sempre esattamente a $p_a$ sul
suo bordo esterno, per definizione. La
Figura~\ref{fig:bellvsplug}c mostra questo, calcolato. I plug reali
sono sempre \emph{troncati} --- tagliati corti e chiusi con una base
piatta --- perché la coda lunga e sottile porta quasi nessuna
pressione ma tutto il peso, e dove tagliare è una genuina variabile
di progetto (Capitolo~\ref{ch:plugcycle}).

\begin{figure}[htbp]
\centering
\includegraphics[width=\textwidth]{fig_bell_vs_plug.pdf}
\caption{Le due famiglie di ugello che questo documento costruisce e
l'unica differenza fisica fra loro. Gli assi sono ovunque posizione
assiale $x$ [m] e raggio $y$ [m]. (a)~La campana: lo scarico è
racchiuso da una parete solida, quindi lo stato di uscita è fissato
dalla geometria, qualunque sia la quota. (b)~Il plug (aerospike): la
spina solida sta sotto e il bordo esterno del getto è una
\emph{frontiera libera} a pressione ambiente. La spina mostrata è un
contorno realmente progettato (dal codice indipendente del
Capitolo~\ref{ch:genotwin}); la frontiera del pennacchio è calcolata
dalla marcia del Capitolo~\ref{ch:plug}; la faccia rossa è la base
dove la spina viene troncata. (c)~Compensazione di quota, calcolata
sullo stesso gas: l'angolo della frontiera del pennacchio $\theta_E$
[gradi] in funzione della pressione ambiente $p_a$ [bar]. Ad
ambiente alto il getto si espande appena e resta aderente alla
spina; al calare dell'ambiente la frontiera ruota verso l'esterno
--- l'ugello che si adatta, senza parti in movimento.}
\label{fig:bellvsplug}
\end{figure}

\section{Il motore a detonazione rotante, e il ciclo}
\label{sec:rdeintro}

Tutto quanto sopra assumeva un solo punto di funzionamento: una
pressione di camera, un ambiente. Un \emph{motore a detonazione
rotante} (RDE) infrange quell'assunzione, ed è la ragione per cui
questo documento ha uno strato di ciclo.

In un RDE la combustione non sta ferma a bruciare in modo stazionario.
Una detonazione --- un'onda di combustione che viaggia più veloce
del suono --- corre di continuo \emph{attorno} a una camera anulare,
percorrendola qualche migliaio di volte al secondo
(Figura~\ref{fig:rdecycle}a). Il propellente fresco viene iniettato
dietro l'onda e consumato al giro successivo. L'attrattiva è
termodinamica: la combustione detonativa alza la pressione mentre
brucia, quindi un RDE promette più lavoro dallo stesso propellente
rispetto a una combustione stazionaria.

La conseguenza per l'ugello è severa. Una sezione fissa dello
scarico non vede uno stato di camera costante; vede il picco di
pressione al passaggio dell'onda e poi il decadimento fino al giro
successivo --- ancora e ancora, migliaia di volte al secondo. Ciò
che viene consegnato all'ugello non è dunque un solo stato di
ristagno ma una \emph{famiglia periodica} di stati,
$(P_0(\xi), T_0(\xi))$, etichettati dalla posizione nel ciclo, la
\emph{fase} $\xi \in [0,1]$ (Figura~\ref{fig:rdecycle}b).

Poiché il gas attraversa l'ugello molto più in fretta di quanto
cambino le condizioni di camera, ogni fase può essere trattata come
un proprio flusso stazionario, e la prestazione dell'intero ciclo è
la media delle prestazioni stazionarie sulle fasi. Quell'ipotesi si
chiama \emph{quasi-stazionaria}; è ciò che rende il problema
trattabile con la macchina di questo documento, è dichiarata
esplicitamente ovunque venga usata, e i suoi limiti sono discussi
onestamente nella Sezione~\ref{sec:spiral}.

La difficoltà di progetto è allora immediata ed è l'intero soggetto
dei Capitoli~\ref{ch:cycle} e~\ref{ch:plugcycle}: ogni fase
vorrebbe un ugello diverso (Figura~\ref{fig:rdecycle}c), e di
hardware se ne costruisce uno solo. Qual è il compromesso giusto ---
ed è semplicemente l'ugello per la condizione media di camera? Per
un tipo di ugello la risposta risulta essere sì, esattamente; per un
altro è decisamente no. Misurare entrambi è ciò che questo documento
fa.

\begin{figure}[htbp]
\centering
\includegraphics[width=\textwidth]{fig_rde_cycle.pdf}
\caption{Perché l'ugello di un RDE è un problema di progetto più
difficile. (a)~L'origine del ciclo: un'onda di detonazione corre
attorno a una camera anulare (gli assi sono coordinate della sezione
di camera $x, y$ [m]); una sezione fissa di ugello la vede passare
una volta per giro, qualche migliaio di volte al secondo. (b)~Ciò
che viene quindi consegnato all'ugello: non un punto di
funzionamento ma una famiglia periodica. Fase di ciclo $\xi$ [--]
sull'asse orizzontale ($\xi = 0$ al passaggio dell'onda), pressione
di camera $P_0$ [bar] sull'asse sinistro (logaritmico, blu) e
temperatura di camera $T_0$ [K] su quello destro (rosso) --- le
cinque fasi discrete effettivamente marciate dai carrier. (c)~Il
problema di progetto che ne risulta: per ogni fase, il rapporto di
espansione $\varepsilon$ [--] che sarebbe perfettamente adattato per
\emph{quella} fase, calcolato con le tabelle di gas del programma.
Si distribuiscono su un fattore tre, e di ugello se ne costruisce
uno solo.}
\label{fig:rdecycle}
\end{figure}

\section{Che cosa significa qui ``ottimizzare'', e che cosa dà un
aggiunto}

Un \emph{ottimizzatore} ha bisogno di tre cose che un solutore di
flusso non possiede. Un \emph{punteggio}: un numero che dica quanto
è buono un progetto --- qui la spinta $J$, costruita nel
Capitolo~\ref{ch:score}. Una \emph{bussola}: la derivata di quel
punteggio rispetto a ogni variabile di progetto, il
\emph{gradiente}, che dice in che direzione muoversi. E un
\emph{certificato}: una condizione che riconosca un ottimo quando lo
si raggiunge, così che ``la ricerca si è fermata'' e ``il progetto è
ottimale'' siano affermazioni distinte.

Il gradiente è il punto in cui la tecnologia del Mattone~1 conta. Se
un progetto ha una sola manopola, il gradiente si può stimare per
\emph{differenze finite} --- si cambia un poco la manopola e si
guarda come cambia $J$. Con un'intera forma di parete come incognita
questo costa una soluzione di flusso in più per ogni grado di
libertà, il che è senza speranza. La \emph{differenziazione
automatica} (AD) evita tutto ciò. Poiché il solutore è scritto in un
linguaggio che registra l'aritmetica che esegue, la regola della
catena può essere applicata all'indietro attraverso l'intero calcolo,
producendo la derivata di $J$ rispetto a \emph{tutti} gli ingressi
in una volta sola, al costo di circa una risoluzione aggiuntiva e
con accuratezza di macchina anziché con quella di un rapporto
incrementale. Questa passata all'indietro è matematicamente lo
stesso oggetto dell'\emph{aggiunto discreto} della teoria classica
del controllo ottimo --- da cui il nome di questo pacchetto di
lavoro --- ma si ottiene gratis anziché derivarlo e programmarlo a
mano.

Questo è il motore. Il resto del documento gli costruisce attorno un
ottimizzatore, un pezzo certificato alla volta.
